\anexo{Transcripción de las entrevistas}\label{anexo:entrevistas}

{\color{nuevo}
Este anexo reúne la transcripción de las cinco entrevistas semiestructuradas realizadas en el marco de la investigación de usuarios, durante agosto de 2026. El análisis de los resultados y las conclusiones derivadas se presentan en el Capítulo~\ref{cap:descripcion}.

Las transcripciones se presentan de forma anonimizada: cada participante se identifica mediante un código (E1, E2, \ldots) y se omite cualquier dato que permita individualizarlo, en línea con el tratamiento de datos personales descrito en ese mismo capítulo. Los participantes prestaron su consentimiento para el registro y la reproducción de sus respuestas con fines académicos.

\subsection*{Consideraciones sobre la transcripción}

\begin{itemize}[leftmargin=2em]
    \item Las intervenciones del entrevistador se indican con \textbf{E:} y las del participante con \textbf{P:}
    \item Se aplicó una \textbf{edición ligera} sobre el registro original: se suprimieron muletillas, titubeos y repeticiones sin valor para el análisis, conservando íntegramente el contenido y el registro coloquial de cada respuesta. No se reformularon ni se completaron ideas del participante.
    \item Se omitió el fragmento inicial de presentación y solicitud de consentimiento, común a las cinco entrevistas, y se sustituyeron los nombres propios y las referencias que permitirían identificar al participante o a su comercio por descripciones genéricas entre corchetes.
    \item Los pasajes entre corchetes corresponden a aclaraciones o a intervenciones no verbales relevantes.
    \item Los puntos suspensivos entre corchetes, \mbox{[\ldots]}, indican la supresión de un fragmento sin valor para el análisis.
\end{itemize}

% ═══════════════════════════════════════════════════════════════════════════
\subsection*{Entrevista 1 (E1)}

\noindent
\begin{tabular}{@{}l@{\hspace{1.5em}}p{10.5cm}@{}}
    \textbf{Perfil} & Remitente particular: revende indumentaria y calzado en marketplaces \\
    \textbf{Fecha} & Agosto de 2026 \\
    \textbf{Modalidad} & Videollamada \\
\end{tabular}

\bigskip

\noindent\textbf{E:} Para empezar, contame un poco de vos y tu relación con los envíos: ¿qué tipo de cosas solés mandar o recibir, y en qué momentos de tu vida aparece esa necesidad?

\noindent\textbf{P:} Mi relación con los envíos consta de mandar y recibir tanto ropa como calzado o accesorios, como pueden ser unos anteojos de sol. Esto surge porque en mi casa somos seis personas y en general nos gusta mucho comprar ropa, pero con el paso del tiempo capaz va quedando chica. Entonces, como está en buen estado, se puede vender para poder volver a comprar más ropa. Y de esa manera no tirar la plata por haber comprado ropa y quizás no haberla usado, o que a alguien le haya quedado chica.

\noindent\textbf{E:} Contame de la última vez que enviaste o recibiste algo, ¿cómo fue el proceso?

\noindent\textbf{P:} La última vez que envié algo fue unos zuecos que me habían regalado mis papás por mi cumpleaños y los compraron afuera; el talle me había quedado chico y entonces los terminé vendiendo. Lo que hice fue publicarlos en Mercado Libre, me llegó una etiqueta, la imprimí, armé mi paquete, la puse, pegué la etiqueta y la llevé al correo para despacharla.

\noindent\textbf{E:} ¿Hubo alguna situación que te generó bronca, dudas o incomodidad? Si es así, ¿podrías contarme qué pasó?

\noindent\textbf{P:} Lo que a veces me puede generar dudas y frustración es que no están bien indicados los lugares en los que se pueden hacer los despachos. O sea, no ponen tan explícitamente en qué lugares son los centros de despacho de los envíos, y eso creo que tanto a mí, que soy joven, como a alguien que capaz es de mayor edad puede generarle confusión.

\noindent\textbf{E:} Cuando tenés que decidir cómo mandar algo, ¿cómo hacés esa elección? ¿Qué aspectos tomás en cuenta?

\noindent\textbf{P:} En general, cuando tengo que enviar algo, a veces puedo elegir el correo, por ejemplo [dos empresas de correo tradicionales], y suelo elegir dependiendo de qué sucursal esté más cerca de mi casa, para no tener que alejarme para despachar el envío.

\noindent\textbf{E:} ¿Alguna vez elegiste a propósito una opción más cara o más lenta? ¿Por qué?

\noindent\textbf{P:} En general elijo la opción estándar. Alguna que otra vez puedo optar por opciones de envío más caras o más lentas por cuestiones de cercanía al centro de despacho o por la confianza que me genera la empresa de logística.

\noindent\textbf{E:} ¿Qué sentís cuando le confiás un paquete a alguien que no conocés? ¿De qué depende tu sensación de seguridad respecto del transportista?

\noindent\textbf{P:} Cuando le doy a alguien que no conozco un paquete, siento mucha incertidumbre. Me ha pasado de tener que realizar envíos por motomensajería, por alguna aplicación, y no tener la certeza de que ese paquete va a llegar a destino y de que la persona que me lo compró lo va a recibir, hasta que efectivamente lo recibe. La verdad, nunca tuve una mala experiencia en la que se perdió el paquete, pero siempre está ese miedo de que puede llegar a no llegar a destino.

\noindent\textbf{E:} Te cuento sobre un proyecto en el que estamos trabajando: una plataforma donde tu envío puede sumarse a un viaje que otra persona ya está haciendo por esa zona, en lugar de contratar un envío exclusivo para vos. En un principio, ¿qué es lo primero que se te viene a la cabeza con respecto a esta idea?

\noindent\textbf{P:} Me parece muy buena la idea, porque de esa manera uno se asegura de que está cubierto por la plataforma, vendría a ser de confianza. Uno le puede dar su paquete y se va a asegurar de que, como ya se está realizando otro recorrido, a la otra persona le va a llegar. Y de esa manera se gana un poco más de confianza y de credibilidad en lo que uno contrata.

\noindent\textbf{E:} Pensándolo un poco más: ¿en qué situaciones te resultaría útil y en cuáles no te animarías? ¿Qué tipo de cosas sí mandás así y cuáles no?

\noindent\textbf{P:} Siendo que yo mando ropa y calzado que ya no se usa en mi casa, esos serían productos que sí enviaría. Capaz, siempre uno tiene miedo de mandar cosas de valor. Si uno vende tecnología o esas cosas, capaz no es la mejor manera de realizar las transacciones. Pero los productos que suelo vender son eso, ropa y calzado, por lo que sí confiaría en mandarlo. Y cuando uno tiene muchos envíos, la verdad, aprovechar los viajes de alguna persona que ya está realizando otros puede ser muy útil.

\noindent\textbf{E:} Para animarte a probar esta plataforma, ¿qué necesitarías que te muestre o te garantice? Y al revés, ¿qué te haría desconfiar o directamente descartar la idea?

\noindent\textbf{P:} Para usar la app, lo que serviría mucho sería un mapa en vivo para ver por dónde está pasando el camión o el transporte que se utilice para realizar los envíos, y los horarios por los que se pasa por cada punto, capaz como ya podemos ver en algunas aplicaciones de envíos tradicionales. Estaría muy bueno también que haya algún sistema en el que se pueda sacar una foto a la persona que lo recibió; o sea, la persona que está haciendo el envío le saca una foto a la persona que lo recibió y hay alguna especie de firma que dé por sentado que la otra persona efectivamente recibió el paquete. Y algo que me haría desconfiar o descartar la idea sería que no haya transparencia, que no haya un aviso del lugar en el que se encuentra el paquete en cada uno de los momentos del trayecto hasta que llegue a destino.

\noindent\textbf{E:} Viéndolo desde el otro lado del proceso: ¿alguna vez se te ocurrió que vos podrías llevar algo de alguien en un viaje que ya ibas a hacer igual? ¿Qué te genera la idea de hacerlo?

\noindent\textbf{P:} Sí, la verdad, se me ocurre que es algo de lo cual podemos sacar ganancia los dos lados. Yo me llevo plata a cambio de dar un servicio que la verdad no me cuesta nada dar, porque de esa manera estoy llevando un producto que probablemente me resulte de camino hacia lo que yo estoy yendo a hacer, que es entregar mi pedido.

\noindent\textbf{E:} Para cerrar, ¿hay algo de tu experiencia con los envíos que no te haya preguntado y creas que nos puede servir?

\noindent\textbf{P:} Yo siento que muchas veces en las aplicaciones de envíos se le da mucha más atención al producto que se está llevando y trayendo que a la experiencia en el transporte en sí. Capaz es importante, dentro de las preguntas que se hacen en las encuestas, agregar cómo fue la experiencia con la empresa de transporte o con quien haya estado encargado de realizar el envío, y de esta manera se pueden mejorar los servicios en general.

\bigskip

% ═══════════════════════════════════════════════════════════════════════════
\subsection*{Entrevista 2 (E2)}

\noindent
\begin{tabular}{@{}l@{\hspace{1.5em}}p{10.5cm}@{}}
    \textbf{Perfil} & Remitente PyME: titular de un comercio gastronómico (franquicia) en el GBA norte \\
    \textbf{Fecha} & Agosto de 2026 \\
    \textbf{Modalidad} & Videollamada \\
\end{tabular}

\bigskip

\noindent\textbf{E:} Contame sobre tu comercio: qué vendés, hace cuánto y cómo es un día normal de trabajo.

\noindent\textbf{P:} Tengo un comercio que es una franquicia gastronómica [de comida japonesa, en el GBA norte]. Hace ya siete años que lo tengo. Un día normal para mí es ir a la mañana a recibir proveedores; después entra el personal y empieza con la producción. Trabajamos solo delivery y \textit{takeaway}, no tengo restaurante al público. Los pedidos empiezan a entrar tipo seis de la tarde hasta las diez, que es el horario en el que más o menos corta la aplicación; trabajo con [una aplicación de delivery]. Esto se da todos los días, abrimos de lunes a lunes; únicamente cambiamos el horario en fechas particulares, como el 24 y el 31 de diciembre, que estamos abiertos solo por la mañana.

\noindent\textbf{E:} Entonces hoy tus productos llegan a manos de tus clientes por delivery o \textit{takeaway}. ¿Podés contarme más sobre cómo se gestionan estos procesos en tu comercio?

\noindent\textbf{P:} Tenemos la aplicación y tenemos una página especial de [la franquicia], que es un carrito en el que te piden lo que quieren; es marcha y sale, o sea, llega el pedido, llega la comanda y se empieza a producir. Y tenemos una o dos motos que reparten, depende del flujo de trabajo, que llevan el pedido a domicilio; o bien hay gente que lo viene a retirar por el local. Y después, como te contaba, trabajamos con la aplicación de delivery: la gente pide por ahí, tenemos media hora para hacer el pedido, llega un motoquero de la aplicación y ellos le entregan directamente al cliente.

\noindent\textbf{E:} Con respecto a eso, nos interesa saber sobre las complicaciones que conllevan estos tipos de envíos. ¿Alguna vez se te complicó algún envío? Si es así, ¿qué pasó y cómo lo terminaste resolviendo?

\noindent\textbf{P:} No. Con respecto a los pedidos, lo que se puede complicar es que se pueden llegar a quejar. Es más, se quejan más para sacar beneficios que por otra cosa; por ejemplo, porque la pieza no es del tamaño adecuado. Pero por suerte no tengo mayores problemas en ese aspecto: de hecho, tengo 4,9 estrellas sobre 5 en la aplicación, y en la página propia también tengo buenas reseñas.

\noindent\textbf{E:} ¿Cómo impacta la disponibilidad de envíos en tu negocio? Contame cómo influye en el precio, en cerrar o perder una venta.

\noindent\textbf{P:} Dentro de los costos operativos hay un sueldo de quien hace envíos con la moto. Con respecto a la aplicación, ellos ya se encargan de pagarle al transportista. Pero nosotros tenemos, como es cadena, un descuento del 5\% sobre la venta de un pedido, que es para la empresa. Generalmente, lo que la aplicación le hace a un comercio normal es un descuento de entre el 13\% y el 15\% del valor del pedido.

\noindent\textbf{E:} Te cuento la idea: que tus productos viajen aprovechando trayectos que otras personas ya hacen por la zona, en vez de un envío exclusivo. Pensando en tu negocio, ¿qué te resulta interesante de la idea y qué te preocupa?

\noindent\textbf{P:} Como vendo alimentos y se requiere un transporte seguro e higiénico, no me sirve directamente para mi producto final. Pero para lo que me serviría es, por ejemplo, si yo pido mercadería en el Mercado Central: podría aprovechar a alguien que realice ya un recorrido cercano y pueda traerme mi mercadería, para aprovechar ese tramo y ahorrar en costos.

\noindent\textbf{E:} Si eso te bajara el costo, pero el que lleva el paquete es un particular y no un cadete o un correo tradicional, ¿cómo lo verías vos?

\noindent\textbf{P:} Si es un particular, o un cadete, o un correo tradicional, estaría perfecto. Si hace su trabajo como corresponde, no tendría problema en que sea quien sea. Es más, lo importante es que las reseñas sean buenas, y con eso me basta y me sobra.

\noindent\textbf{E:} Para usarlo semana a semana y confiarle tus envíos, ¿qué no podría faltar? ¿Y qué sería una señal de que es mejor seguir como estás hoy?

\noindent\textbf{P:} La prioridad en ese caso sería la transparencia: un mapa que muestre la ubicación del transportista en cada momento y que los perfiles estén completos para identificar fácilmente a la persona. No usaría la aplicación si, por ejemplo, no puedo darle seguimiento a mi envío en tiempo real.

\noindent\textbf{E:} ¿Hay algo de la logística de tu negocio que no hayamos charlado y te parezca importante?

\noindent\textbf{P:} Por el momento, no.

\bigskip

% ═══════════════════════════════════════════════════════════════════════════
\subsection*{Entrevista 3 (E3)}

\noindent
\begin{tabular}{@{}l@{\hspace{1.5em}}p{10.5cm}@{}}
    \textbf{Perfil} & Transporte de carga: conductor de camión, en relación de dependencia con una empresa de reparto \\
    \textbf{Fecha} & Agosto de 2026 \\
    \textbf{Modalidad} & Videollamada \\
\end{tabular}

\bigskip

\noindent\textbf{E:} Para ubicarme, contame un poco: ¿el transporte es algo que hacés como trabajo, o sos más de moverte por tus cosas y capaz podrías aprovechar esos viajes?

\noindent\textbf{P:} El transporte lo hago como trabajo, sí, como medio de vida.

\noindent\textbf{E:} Contame cómo es un día o una semana típica tuya moviéndote: ¿qué zonas recorrés, en qué te movés, cómo son esos recorridos? ¿Cómo conseguís hoy los viajes o los clientes?

\noindent\textbf{P:} Me muevo por donde requieran el reparto, por donde estén los clientes. No tengo un recorrido fijo y me muevo en un camión. Y esos recorridos, o esos clientes, los maneja una empresa, la cual tiene gente que levanta pedidos y nosotros hacemos los repartos.

\noindent\textbf{E:} Te cuento la idea: aprovechar un viaje que ya vas a hacer igual para llevar una encomienda que te quede sobre el recorrido, a cambio de una compensación. ¿Qué pensás?

\noindent\textbf{P:} Si queda sobre el recorrido y no varía nada, se puede hacer por el mismo precio del flete, siempre y cuando quede en el mismo recorrido.

\noindent\textbf{E:} ¿En qué condiciones te sumarías y en cuáles ni lo pensarías? Contame qué tendría que pasar para que digas «esto me sirve».

\noindent\textbf{P:} Las condiciones para decir «esto no me sirve» o «esto me sirve» serían cuando tenés un recorrido y quisieran que salgas de ese recorrido y te desvíes muchísimo para entregar un pedido, el cual no te lo van a pagar. Ese desgaste y esos gastos los tenés que cubrir vos y no te los pagan.

\noindent\textbf{E:} Cuando pensás en llevar algo de alguien que no conocés, ¿qué es lo que más te haría ruido?

\noindent\textbf{P:} Lo que más me haría ruido es cuando no importa el valor, llevarlo igual, y no conocés al destinatario ni a la persona que lo envía, y no tiene ninguna factura en blanco como para que acredite lo que realmente está llevando. Y si se daña o se pierde, se tiene que hacer cargo el seguro de la mercadería. La mercadería, por lo general, siempre sale asegurada.

\noindent\textbf{E:} Antes de aceptar un envío, ¿qué necesitás saber sí o sí para decidir si lo tomás?

\noindent\textbf{P:} Antes de aceptar un envío, lo que sí o sí hay que acordar es el costo del flete o el costo del traslado, más todo lo que ello implique, y que sí o sí acredite con una facturación en blanco qué estoy llevando, para no encontrarme con sorpresas.

\noindent\textbf{E:} ¿Cómo te das cuenta de que un pago es justo? Contame qué tenés en cuenta para poner un precio.

\noindent\textbf{P:} Para saber si el pago es justo, lo que hay que tener en cuenta son todos los gastos de logística que tenés para transportar la mercadería que te designen: desgaste del vehículo, precio de combustible, peajes, en el caso de si necesitás algún ayudante. Todas esas cosas hay que tenerlas en cuenta al momento de poner un costo para poder hacer el viaje.

\noindent\textbf{E:} ¿Qué tanto te importa que la plataforma no te penalice por no aceptar un pedido o por desconectarte un tiempo? Contame.

\noindent\textbf{P:} Es importante que no te penalicen o que no te cobren por entrar y salir, porque hay días que por ahí podés tener viajes por otro lado o para otra empresa, y que se te remunere mejor, y no hacer viaje para la app ese día.

\noindent\textbf{E:} ¿Hay algo de tu experiencia transportando que no te haya preguntado y te parezca importante compartir?

\noindent\textbf{P:} Sí: que las apps tienen que ser específicas con todo lo que se le va a cobrar al cliente al momento de darte un viaje de transporte o de lo que fuere. Tiene que ser específico de que el viaje va a salir X cantidad de plata, tendrá gastos fijos, como lo son los peajes. Muchas apps no tienen en cuenta los peajes, que hoy por hoy también es un gasto importante. Y muchas apps no tienen en cuenta lugares en los que hay que estacionar, que hay que pagar estacionamiento para poder descargar, o poder llevar gente, o poder esperar. Y todo eso tiene que estar incluido. La app tendría que tener un registro de dónde se va a parar el vehículo, si tiene que pagar estacionamiento medido, si tiene peajes y todos los gastos, para poder pasarlos al cliente.

\bigskip

% ═══════════════════════════════════════════════════════════════════════════
\subsection*{Entrevista 4 (E4)}

\noindent
\begin{tabular}{@{}l@{\hspace{1.5em}}p{10.5cm}@{}}
    \textbf{Perfil} & Destinataria particular: compra en línea con frecuencia y proyecta comenzar a vender \\
    \textbf{Fecha} & Agosto de 2026 \\
    \textbf{Modalidad} & Videollamada \\
\end{tabular}

\bigskip

\noindent\textbf{E:} Para empezar, contame un poco de vos y tu relación con los envíos: ¿qué tipo de cosas solés mandar o recibir, y en qué momentos de tu vida aparece esa necesidad?

\noindent\textbf{P:} Por lo general solo recibo; hasta el momento no tuve la necesidad de enviar, pero sí recibo pedidos que hago a través de páginas de internet. Más que nada recibo compras de elementos para la casa o, en oportunidades, productos que no consigo fácilmente en los locales cercanos, como cápsulas de café.

\noindent\textbf{E:} Contame de la última vez que enviaste o recibiste algo, ¿cómo fue el proceso?

\noindent\textbf{P:} La última vez que recibí algo fue un elemento para la cocina. Decidí comprarlo a través de la página oficial de la marca que lo vende, y solicité que me envíen el paquete a una sucursal de una empresa de logística que queda cerca de mi casa, ya que el producto era costoso y no quería arriesgarme a que sufriera daños y que no estuvieran cubiertos.

\noindent\textbf{E:} ¿Hubo alguna situación que te generó bronca, dudas o incomodidad? Si es así, ¿podrías contarme qué pasó?

\noindent\textbf{P:} Una vez tuve una situación compleja con una compra hecha a través de un marketplace, a la cual le robaron al transportista cuando se trasladaba en su moto. El vendedor no era fácil de contactar y, una vez que logré hablarle, quería que aceptara que el pedido quedara marcado como «entregado» mientras me enviaba el reemplazo, a lo que me negué.

\noindent\textbf{E:} Cuando tenés que decidir cómo recibir algo, ¿cómo hacés esa elección? ¿Qué aspectos tomás en cuenta?

\noindent\textbf{P:} Suelo comprar a través de las páginas oficiales de internet, entonces evalúo qué opciones de envío tienen disponibles. Siempre trato de optar por empresas de logística conocidas y que tengan sucursales cercanas a mi casa, para evitar la espera en mi casa hasta que llegue el pedido.

\noindent\textbf{E:} ¿Alguna vez elegiste a propósito una opción más cara o más lenta? ¿Por qué?

\noindent\textbf{P:} La opción del envío a sucursal, si bien es más barata, suele ser más lenta, pero me deja más tranquila, porque no tengo que estar en mi casa esperando durante muchas horas un día, o durante varios días, para recibir el pedido, ya que suelen avisar sobre el momento de entrega.

\noindent\textbf{E:} ¿Qué sentís cuando le confiás un paquete a alguien que no conocés? ¿De qué depende tu sensación de seguridad respecto del transportista?

\noindent\textbf{P:} Honestamente, no me deja muy tranquila. Me da más seguridad cuando el envío cuenta con seguimiento a través de una plataforma.

\noindent\textbf{E:} Te cuento sobre un proyecto en el que estamos trabajando: una plataforma donde tu envío puede sumarse a un viaje que otra persona ya está haciendo por esa zona, en lugar de contratar un envío exclusivo para vos. En un principio, ¿qué es lo primero que se te viene a la cabeza con respecto a esta idea?

\noindent\textbf{P:} Me gusta mucho la idea, me parece que sería una gran opción para reducir gastos y generar una red local de logística.

\noindent\textbf{E:} Pensándolo un poco más: ¿en qué situaciones te resultaría útil y en cuáles no te animarías? ¿Qué tipo de cosas sí mandás así y cuáles no?

\noindent\textbf{P:} Me parece que sería útil para enviar y recibir paquetes entre municipios cercanos, sin tener que hacer el trayecto en el vehículo propio. No enviaría cosas de gran valor o importantes; pero, por ejemplo, quiero empezar a vender cosas que no uso, como ropa y libros, y esos elementos podrían ser enviados sin problema.

\noindent\textbf{E:} Para animarte a probar esta plataforma, ¿qué necesitarías que te muestre o te garantice? Y al revés, ¿qué te haría desconfiar o directamente descartar la idea?

\noindent\textbf{P:} Sí o sí debería tener un seguimiento en tiempo real del envío y un perfil con datos para identificar al transportista. Teniendo en cuenta que los transportistas serían personas que lo hacen de forma independiente, debe haber una manera de comunicarse con la persona en caso de haber un problema. No tener esa información me haría dudar de la confiabilidad del transportista.

\noindent\textbf{E:} Viéndolo desde el otro lado del proceso: ¿alguna vez se te ocurrió que vos podrías llevar algo de alguien en un viaje que ya ibas a hacer igual? ¿Qué te genera la idea de hacerlo?

\noindent\textbf{P:} Me parece una propuesta interesante, no la había pensado. Hace poco me retiré de mi trabajo y tengo más tiempo para hacer actividades. Podría aprovechar esos trayectos, o incluso otros momentos libres, para generar un ingreso extra.

\noindent\textbf{E:} Para cerrar, ¿hay algo de tu experiencia con los envíos que no te haya preguntado y creas que nos puede servir?

\noindent\textbf{P:} Sí. El otro día me di cuenta de cuánto contamina el rubro de la logística, entre viajes y embalaje. Me parece que la idea de sacarle el mayor provecho a los viajes es una buena oportunidad de mejorar esta situación.

\bigskip

% ═══════════════════════════════════════════════════════════════════════════
\subsection*{Entrevista 5 (E5)}

\noindent
\begin{tabular}{@{}l@{\hspace{1.5em}}p{10.5cm}@{}}
    \textbf{Perfil} & Transportista con recorrido habitual: distribuidor independiente con camioneta, zona sur y oeste del AMBA \\
    \textbf{Fecha} & Agosto de 2026 \\
    \textbf{Modalidad} & Videollamada \\
\end{tabular}

\bigskip

\noindent\textbf{E:} Contame un poco: ¿el transporte es algo que hacés como trabajo, o sos más de moverte por tus cosas y capaz podrías aprovechar esos viajes?

\noindent\textbf{P:} Mi trabajo es ser el nexo entre los que manufacturan un producto y quienes lo comercializan; yo lo distribuyo a distintos comercios.

\noindent\textbf{E:} Contame cómo es un día o una semana típica tuya moviéndote: ¿qué zonas recorrés, en qué te movés, cómo son esos recorridos? ¿Cómo conseguís hoy los viajes o los clientes?

\noindent\textbf{P:} Por lo general me muevo en la zona sur del AMBA, aunque también trabajo con algunos clientes en la zona oeste. Uso mi vehículo particular, que es una camioneta. Me dedico a un nicho muy específico del mercado, entonces tanto a los productores como a los clientes que comercializan los busqué yo, y les ofrezco mi servicio de logística.

\noindent\textbf{E:} Entonces entiendo que trabajás de forma independiente, ¿es así?

\noindent\textbf{P:} Exactamente, por eso yo suelo buscar los clientes en lugar de ser al revés.

\noindent\textbf{E:} Te cuento la idea: aprovechar un viaje que ya vas a hacer igual para llevar una encomienda que te quede sobre el recorrido, a cambio de una compensación. ¿Qué pensás?

\noindent\textbf{P:} Me parece una buena idea. Se me ocurre que podría usarlo los días que salgo a hacer los repartos, ya que recorro varios lugares y hago distancias considerables.

\noindent\textbf{E:} ¿En qué condiciones te sumarías y en cuáles ni lo pensarías? Contame qué tendría que tener la plataforma para que te haga sumarte.

\noindent\textbf{P:} Lo haría solo si puedo utilizar mi recorrido preestablecido, o usarlo con pequeños desvíos. Claramente no lo usaría si me requiere desviarme y perder mucho tiempo en la entrega, porque tengo que visitar varios clientes en un día y eso es prioritario.

\noindent\textbf{E:} Cuando pensás en llevar algo de alguien que no conocés, ¿qué es lo que más te haría ruido?

\noindent\textbf{P:} Principalmente, el hecho de la responsabilidad en caso de que lo que se lleva sufra algún daño. Por otra parte, si no hay un registro o una garantía, la persona que debe recibir el paquete puede hacer acusaciones falsas de robo, o quejas por daños que no existen, para conseguir un beneficio económico.

\noindent\textbf{E:} Antes de aceptar un envío, ¿qué necesitás saber sí o sí para decidir si lo tomás?

\noindent\textbf{P:} Debería ver si queda bien establecido con mi recorrido, qué es lo que se quiere transportar, para saber si tengo lugar en el vehículo, y el precio, obviamente.

\noindent\textbf{E:} ¿Qué tendrías en cuenta para calcular un precio que te convenga? Contame qué tenés en cuenta hoy para poner un precio.

\noindent\textbf{P:} Hoy en día el cálculo lo hago para obtener un margen de ganancia entre el precio al que me vendió el productor y el precio al que se lo vendo a los comerciantes; a eso se le suman los gastos de transporte. Teniendo en cuenta que los envíos, en este caso, serían un extra sobre mi trabajo habitual, buscaría cubrir con ellos los gastos de nafta y de mantenimiento de la camioneta.

\noindent\textbf{E:} ¿Qué tanto te importa que la plataforma no te penalice por no aceptar un pedido o por desconectarte un tiempo?

\noindent\textbf{P:} En mi caso me importa muchísimo, porque lo haría como un extra, solo los días que salgo a repartir a mis clientes y solo si me resulta conveniente. No sería un trabajo fijo.

\noindent\textbf{E:} ¿Hay algo de tu experiencia transportando que no te haya preguntado y te parezca importante compartir?

\noindent\textbf{P:} Por el momento no se me viene nada más a la cabeza.
}
